\section{Completion}
\label{section-completion}

\noindent
Suppose that $R$ is a ring and $I$ is an ideal.
We define the {\it completion of $R$ with respect to $I$}
to be the limit
$$
R^\wedge = \lim_n R/I^n.
$$
An element of $R^\wedge$ is given by a sequence
of elements $f_n \in R/I^n$ such that $f_n \equiv f_{n + 1} \bmod I^n$
for all $n$. We will view $R^\wedge$ as an $R$-algebra.
Similarly, if $M$ is an $R$-module then we define the
{\it completion of $M$ with respect to $I$}
to be the limit
$$
M^\wedge = \lim_n M/I^nM.
$$
An element of $M^\wedge$ is given by a sequence of
elements $m_n \in M/I^nM$ such that $m_n \equiv m_{n + 1} \bmod I^nM$
for all $n$. We will view $M^\wedge$ as an $R^\wedge$-module.
From this description it is clear that there
are always canonical maps
$$
M \longrightarrow M^\wedge
\quad\text{and}\quad
M \otimes_R R^\wedge \longrightarrow M^\wedge.
$$
Moreover, given a map $\varphi : M \to N$ of modules we get an induced
map $\varphi^\wedge : M^\wedge \to N^\wedge$ on completions making the
diagram
$$
\xymatrix{
M \ar[r] \ar[d] & N \ar[d] \\
M^\wedge \ar[r] & N^\wedge
}
$$
commute. In general completion is not an exact functor, see
Examples, Section \ref{examples-section-completion-not-exact}.
Here are some initial positive results.

\begin{lemma}
\label{lemma-completion-generalities}
Let $R$ be a ring. Let $I \subset R$ be an ideal.
Let $\varphi : M \to N$ be a map of $R$-modules.
\begin{enumerate}
\item If $M/IM \to N/IN$ is surjective, then $M^\wedge \to N^\wedge$
is surjective.
\item If $M \to N$ is surjective, then $M^\wedge \to N^\wedge$ is surjective.
\item If $0 \to K \to M \to N \to 0$ is a short exact sequence of
$R$-modules and $N$ is flat, then
$0 \to K^\wedge \to M^\wedge \to N^\wedge \to 0$ is a short exact sequence.
\item The map $M \otimes_R R^\wedge \to M^\wedge$ is
surjective for any finite $R$-module $M$.
\end{enumerate}
\end{lemma}

\begin{proof}
Assume $M/IM \to N/IN$ is surjective. Then the map $M/I^nM \to N/I^nN$
is surjective for each $n \geq 1$ by Nakayama's lemma. More precisely,
apply Lemma \ref{lemma-NAK} part (11) to the
map $M/I^nM \to N/I^nN$ over the ring $R/I^n$ and the nilpotent
ideal $I/I^n$ to see this. Set $K_n = \{x \in M \mid \varphi(x) \in I^nN\}$.
Thus we get short exact sequences
$$
0 \to K_n/I^nM \to M/I^nM \to N/I^nN \to 0
$$
We claim that the canonical map $K_{n + 1}/I^{n + 1}M \to K_n/I^nM$
is surjective. Namely, if $x \in K_n$ write
$\varphi(x) = \sum z_j n_j$ with $z_j \in I^n$, $n_j \in N$.
By assumption we can write $n_j = \varphi(m_j) + \sum z_{jk}n_{jk}$
with $m_j \in M$, $z_{jk} \in I$ and $n_{jk} \in N$. Hence
$$
\varphi(x - \sum z_j m_j) = \sum z_jz_{jk} n_{jk}.
$$
This means that $x' = x - \sum z_j m_j \in K_{n + 1}$ maps
to $x \bmod I^nM$ which proves the claim. Now we may apply
Lemma \ref{lemma-Mittag-Leffler}
to the inverse system of short exact sequences above to see (1).
Part (2) is a special case of (1).
If the assumptions of (3) hold, then for each $n$ the sequence
$$
0 \to K/I^nK \to M/I^nM \to N/I^nN \to 0
$$
is short exact by
Lemma \ref{lemma-flat-tor-zero}.
Hence we can directly apply
Lemma \ref{lemma-Mittag-Leffler}
to conclude (3) is true.
To see (4) choose generators $x_i \in M$, $i = 1, \ldots, n$.
Then the map $R^{\oplus n} \to M$, $(a_1, \ldots, a_n) \mapsto \sum a_ix_i$
is surjective. Hence by (2) we see
$(R^\wedge)^{\oplus n} \to M^\wedge$, $(a_1, \ldots, a_n) \mapsto \sum a_ix_i$
is surjective. Assertion (4) follows from this.
\end{proof}

\begin{definition}
\label{definition-complete}
Let $R$ be a ring. Let $I \subset R$ be an ideal.
Let $M$ be an $R$-module. We say $M$ is {\it $I$-adically complete}
if the map
$$
M \longrightarrow M^\wedge = \lim_n M/I^nM
$$
is an isomorphism\footnote{This includes the condition that
$\bigcap I^nM = 0$.}. We say $R$ is {\it $I$-adically complete}
if $R$ is $I$-adically complete as an $R$-module.
\end{definition}

\noindent
It is not true that the completion of an $R$-module $M$ with respect
to $I$ is $I$-adically complete. For an example see
Examples, Section \ref{examples-section-noncomplete-completion}.
If the ideal is finitely generated, then the completion is complete.

\begin{lemma}
\label{lemma-hathat-finitely-generated}
\begin{reference}
\cite[Theorem 15]{Matlis}. The slick proof given here is from
an email of Bjorn Poonen dated Nov 5, 2016.
\end{reference}
Let $R$ be a ring. Let $I$ be a finitely generated ideal of $R$.
Let $M$ be an $R$-module. Then
\begin{enumerate}
\item the completion $M^\wedge$ is $I$-adically complete, and
\item $I^nM^\wedge = \Ker(M^\wedge \to M/I^nM) = (I^nM)^\wedge$ for all
$n \geq 1$.
\end{enumerate}
In particular $R^\wedge$ is $I$-adically complete,
$I^nR^\wedge = (I^n)^\wedge$, and
$R^\wedge/I^nR^\wedge = R/I^n$.
\end{lemma}

\begin{proof}
Since $I$ is finitely generated,
$I^n$ is finitely generated, say by $f_1, \ldots, f_r$.
Applying Lemma \ref{lemma-completion-generalities} part (2)
to the surjection $(f_1, \ldots, f_r) : M^{\oplus r} \to I^n M$
yields a surjection
$$
(M^\wedge)^{\oplus r} \xrightarrow{(f_1, \ldots, f_r)} (I^n M)^\wedge =
\lim_{m \geq n} I^n M/I^m M = \Ker(M^\wedge \to M/I^n M).
$$
On the other hand, the image of
$(f_1, \ldots, f_r) : (M^\wedge)^{\oplus r} \to M^\wedge$
is $I^n M^\wedge$.
Thus $M^\wedge / I^n M^\wedge \simeq M/I^n M$.
Taking inverse limits yields $(M^\wedge)^\wedge \simeq M^\wedge$;
that is, $M^\wedge$ is $I$-adically complete.
\end{proof}

\begin{lemma}
\label{lemma-completion-differ-by-torsion}
Let $R$ be a ring. Let $I \subset R$ be an ideal. Let
$0 \to M \to N \to Q \to 0$ be an exact sequence of
$R$-modules such that $Q$ is annihilated by a power of $I$.
Then completion produces an exact sequence
$0 \to M^\wedge \to N^\wedge \to Q \to 0$.
\end{lemma}

\begin{proof}
Say $I^c Q = 0$. Then $Q/I^nQ = Q$ for $n \geq c$.
On the other hand, it is clear that
$I^nM \subset M \cap I^nN \subset I^{n - c}M$ for $n \geq c$.
Thus $M^\wedge = \lim M/(M \cap I^n N)$. Apply Lemma \ref{lemma-Mittag-Leffler}
to the system of exact sequences
$$
0 \to M/(M \cap I^n N) \to N/I^n N \to Q \to 0
$$
for $n \geq c$ to conclude.
\end{proof}

\begin{lemma}
\label{lemma-hathat}
\begin{reference}
Taken from an unpublished note of Lenstra and de Smit.
\end{reference}
Let $R$ be a ring. Let $I \subset R$ be an ideal. Let $M$ be an $R$-module.
Denote $K_n = \Ker(M^\wedge \to M/I^nM)$. Then $M^\wedge$ is $I$-adically
complete if and only if $K_n$ is equal to $I^nM^\wedge$ for all $n \geq 1$.
\end{lemma}

\begin{proof}
The module $I^n M^\wedge$ is contained in $K_n$.
Thus for each $n \geq 1$ there is a canonical exact sequence
$$
0 \to K_n/I^nM^\wedge \to M^\wedge/I^nM^\wedge \to M/I^nM \to 0.
$$
As $I^nM^\wedge$ maps onto $I^nM/I^{n + 1}M$ we see that
$K_{n + 1} + I^n M^\wedge = K_n$. Thus the inverse system
$\{K_n/I^n M^\wedge\}_{n \geq 1}$ has surjective transition maps.
By
Lemma \ref{lemma-Mittag-Leffler}
we see that there is a short exact sequence
$$
0 \to
\lim_n K_n/I^n M^\wedge \to
(M^\wedge)^\wedge \to
M^\wedge \to 0
$$
Hence $M^\wedge$ is complete if and only if $K_n/I^n M^\wedge = 0$
for all $n \geq 1$.
\end{proof}

\begin{lemma}
\label{lemma-radical-completion}
Let $R$ be a ring, let $I \subset R$ be an ideal, and let
$R^\wedge = \lim R/I^n$.
\begin{enumerate}
\item any element of $R^\wedge$ which maps to a unit of $R/I$ is a unit,
\item any element of $1 + I$ maps to an invertible element of $R^\wedge$,
\item any element of $1 + IR^\wedge$ is invertible in $R^\wedge$, and
\item the ideals $IR^\wedge$ and $\Ker(R^\wedge \to R/I)$ are contained
in the Jacobson radical of $R^\wedge$.
\end{enumerate}
\end{lemma}

\begin{proof}
Let $x \in R^\wedge$ map to a unit $x_1$ in $R/I$.
Then $x$ maps to a unit $x_n$ in $R/I^n$ for every $n$ by
Lemma \ref{lemma-locally-nilpotent-unit}.
Hence $y = (x_n^{-1}) \in \lim R/I^n = R^\wedge$ is an inverse to $x$.
Parts (2) and (3) follow immediately from (1).
Part (4) follows from (1) and Lemma \ref{lemma-contained-in-radical}.
\end{proof}

\begin{lemma}
\label{lemma-when-surjective-to-completion}
Let $A$ be a ring. Let $I = (f_1, \ldots, f_r)$ be a finitely
generated ideal. If $M \to \lim M/f_i^nM$ is surjective for
each $i$, then $M \to \lim M/I^nM$ is surjective.
\end{lemma}

\begin{proof}
Note that $\lim M/I^nM = \lim M/(f_1^n, \ldots, f_r^n)M$ as
$I^n \supset (f_1^n, \ldots, f_r^n) \supset I^{rn}$.
An element $\xi$ of $\lim M/(f_1^n, \ldots, f_r^n)M$ can be symbolically
written as
$$
\xi = \sum\nolimits_{n \geq 0} \sum\nolimits_i f_i^n x_{n, i}
$$
with $x_{n, i} \in M$. If $M \to \lim M/f_i^nM$ is surjective, then there is
an $x_i \in M$ mapping to $\sum x_{n, i} f_i^n$ in $\lim M/f_i^nM$.
Then $x = \sum x_i$ maps to $\xi$ in $\lim M/I^nM$.
\end{proof}

\begin{lemma}
\label{lemma-complete-by-sub}
Let $A$ be a ring. Let $I \subset J \subset A$ be ideals.
If $M$ is $J$-adically complete and $I$ is finitely generated, then
$M$ is $I$-adically complete.
\end{lemma}

\begin{proof}
Assume $M$ is $J$-adically complete and $I$ is finitely generated.
We have $\bigcap I^nM = 0$ because $\bigcap J^nM = 0$. By
Lemma \ref{lemma-when-surjective-to-completion}
it suffices to prove the surjectivity of $M \to \lim M/I^nM$ in case
$I$ is generated by a single element. Say $I = (f)$.
Let $x_n \in M$ with $x_{n + 1} - x_n \in f^nM$. We have to show there exists
an $x \in M$ such that $x_n - x \in f^nM$ for all $n$.
As $x_{n + 1} - x_n \in J^nM$ and as $M$ is $J$-adically complete,
there exists an element $x \in M$ such that $x_n - x \in J^nM$.
Replacing $x_n$ by $x_n - x$ we may assume that $x_n \in J^nM$.
To finish the proof we will show that this implies $x_n \in I^nM$.
Namely, write $x_n - x_{n + 1} = f^nz_n$.
Then
$$
x_n = f^n(z_n + fz_{n + 1} + f^2z_{n + 2} + \ldots)
$$
The sum $z_n + fz_{n + 1} + f^2z_{n + 2} + \ldots$ converges in $M$
as $f^c \in J^c$. The sum $f^n(z_n + fz_{n + 1} + f^2z_{n + 2} + \ldots)$
converges in $M$ to $x_n$ because
the partial sums equal $x_n - x_{n + c}$ and $x_{n + c} \in J^{n + c}M$.
\end{proof}

\begin{lemma}
\label{lemma-change-ideal-completion}
Let $R$ be a ring.
Let $I$, $J$ be ideals of $R$.
Assume there exist integers $c, d > 0$ such that
$I^c \subset J$ and $J^d \subset I$.
Then completion with respect to $I$ agrees with completion
with respect to $J$ for any $R$-module.
In particular an $R$-module $M$ is $I$-adically complete
if and only if it is $J$-adically complete.
\end{lemma}

\begin{proof}
Consider the system of maps
$M/I^nM \to M/J^{\lfloor n/c \rfloor}M$ and
the system of maps $M/J^mM \to M/I^{\lfloor m/d \rfloor}M$
to get mutually inverse maps between the completions.
\end{proof}

\begin{lemma}
\label{lemma-quotient-complete}
Let $R$ be a ring. Let $I$ be an ideal of $R$.
Let $M$ be an $I$-adically complete $R$-module,
and let $K \subset M$ be an $R$-submodule.
The following are equivalent
\begin{enumerate}
\item $K = \bigcap (K + I^nM)$ and
\item $M/K$ is $I$-adically complete.
\end{enumerate}
\end{lemma}

\begin{proof}
Set $N = M/K$. By
Lemma \ref{lemma-completion-generalities}
the map $M = M^\wedge \to N^\wedge$ is surjective.
Hence $N \to N^\wedge$ is surjective. It is easy to see that the
kernel of $N \to N^\wedge$ is the module $\bigcap (K + I^nM) / K$.
\end{proof}

\begin{lemma}
\label{lemma-when-finite-module-complete-over-complete-ring}
Let $R$ be a ring. Let $I$ be an ideal of $R$.
Let $M$ be an $R$-module.
If (a) $R$ is $I$-adically complete, (b) $M$ is a finite $R$-module,
and (c) $\bigcap I^nM = (0)$, then $M$ is $I$-adically complete.
\end{lemma}

\begin{proof}
By Lemma \ref{lemma-completion-generalities}
the map $M = M \otimes_R R = M \otimes_R R^\wedge \to M^\wedge$
is surjective. The kernel of this map is $\bigcap I^nM$ hence zero
by assumption. Hence $M \cong M^\wedge$ and $M$ is complete.
\end{proof}

\begin{lemma}
\label{lemma-finite-over-complete-ring}
Let $R$ be a ring. Let $I \subset R$ be an ideal. Let $M$ be an $R$-module.
Assume
\begin{enumerate}
\item $R$ is $I$-adically complete,
\item $\bigcap_{n \geq 1} I^nM = (0)$, and
\item $M/IM$ is a finite $R/I$-module.
\end{enumerate}
Then $M$ is a finite $R$-module.
\end{lemma}

\begin{proof}
Let $x_1, \ldots, x_n \in M$ be elements whose images in $M/IM$ generate
$M/IM$ as a $R/I$-module. Denote $M' \subset M$ the $R$-submodule
generated by $x_1, \ldots, x_n$. By Lemma \ref{lemma-completion-generalities}
the map $(M')^\wedge \to M^\wedge$ is surjective.
Since $\bigcap I^nM = 0$ we see in particular that $\bigcap I^nM' = (0)$.
Hence by Lemma \ref{lemma-when-finite-module-complete-over-complete-ring}
we see that $M'$ is complete, and we conclude that $M' \to M^\wedge$
is surjective. Finally, the kernel of $M \to M^\wedge$ is
zero since it is equal to $\bigcap I^nM = (0)$.
Hence we conclude that $M \cong M' \cong M^\wedge$
is finitely generated.
\end{proof}









